% Calibration corpus for scholar-verifier step 7.6 (claim <-> own evidence).
% Each sentence carries an id comment so a grader can report unambiguously.
% Ground truth lives in ground_truth.json and is NOT to be shown to a grader.

\section{Method}

% S01
We train each policy for 300 epochs with a batch size of 4096.

\section{Experiments}

\subsection{Setup}

% S02
Each configuration is evaluated over 200 episodes with five random seeds.
% S03
Prior work reports a 12\% drop under the same distribution shift~\cite{smith2024}.

\subsection{Main results}

% S04
Table~\ref{tab:main} reports success rate on all six terrains.
% S05
Our method leads on five of the six.
% S06
Our controller is substantially more robust than the baselines.
% S07
Table~\ref{tab:main} demonstrates that our method is universally superior.
% S08
The 1.4-point difference in Table~\ref{tab:main} suggests that the auxiliary loss helps, though the effect is within run-to-run variance.

\subsection{Ablations}

% S09
Success rate rises from 41.2\% to 58.7\% when the auxiliary loss is enabled.
% S10
Performance improves considerably once the auxiliary loss is enabled.
% S11
In this single ablation, success rate rises from 41\% to 68\%, which proves that curriculum scheduling is necessary for convergence.
% S12
The 3.2-point gain in Table~\ref{tab:ablation} establishes the general effectiveness of the proposed regulariser.

\subsection{Transfer}

% S13
Figure~\ref{fig:shift} confirms that the policy is robust to any distribution shift.
% S14
The proposed method generalizes well to unseen terrain.
% S15
Across all three benchmarks (Tables~\ref{tab:main}, \ref{tab:ablation}, and~\ref{tab:transfer}), the method matches or exceeds the strongest baseline, which demonstrates that the gain is not benchmark-specific.
% S16
Our approach outperforms all prior work in this setting.

\subsection{Limitations}

% S17
We did not evaluate on hardware, so these numbers may not transfer to the physical platform.
